\section{Očekivanja i stvarnost: od 2003. do 2026. godine}
\label{sec:retrospektiva}

Izbor Rijndaela za Advanced Encryption Standard (AES) završen je 2000.
godine, a standard FIPS~197 objavljen je 2001. godine. Zbog toga se
dokumentovana osnova za rekonstrukciju početnih očekivanja nalazi pretežno
u izvorima iz perioda 1997--2001, dok period od približno 2003. do 2026.
predstavlja vremenski okvir u kome se može posmatrati šta se sa tim
očekivanjima dogodilo u praksi. NIST-ov izveštaj o razvoju AES-a pokazuje
da izbor nije bio zasnovan na jednoj metričkoj funkciji ili numeričkom
rangiranju, već na kombinaciji bezbednosti, troška implementacije,
softverskih i hardverskih karakteristika i fleksibilnosti
\cite{nechvatal2001aes}. FIPS~197 je zatim standardizovao AES-128,
AES-192 i AES-256, sa zajedničkom veličinom bloka od 128 bita
\cite{nist_fips197_2001,nist_fips197_2023}.

Retrospektiva zato ne pokušava da proglasi originalna predviđanja
``tačnim'' ili ``pogrešnim'' u apsolutnom smislu. Umesto toga, porede se
dokumentovana očekivanja i kriterijumi iz procesa standardizacije sa
kasnijim proverljivim ishodima. Posebno se razdvajaju svojstva same
kriptografske primitive, svojstva konkretne implementacije i svojstva
protokola koji primitivu koristi. Isto razdvajanje bilo je ključno i u
eksperimentalnom delu ovog rada: performanse softverske AES implementacije,
AES-NI putanje, CUDA kernela i kompletne CPU--GPU putanje ne predstavljaju
istu veličinu, kao što ni otpornost AES primitive ne garantuje ispravnu
upotrebu režima rada ili protokola.

\subsection{Rekonstrukcija očekivanja}

NIST je proces razvoja AES-a započeo 1997. godine sa ciljem izbora
simetričnog algoritma koji bi zamenio DES i ostao upotrebljiv tokom
dugog vremenskog perioda. Zahtevan je najmanje 128-bitni blok i podrška
za ključeve od 128, 192 i 256 bita, a kandidati su ocenjivani po
bezbednosti, trošku i karakteristikama algoritma i implementacije
\cite{nechvatal2001aes}. Bezbednost je imala najveću težinu, ali NIST
nije posmatrao samo broj rundi ili veličinu ključa: razmatrani su poznati
napadi, bezbednosna margina, jednostavnost strukture, statističko
ponašanje, softverske i hardverske performanse, memorijski zahtevi i
pogodnost za različite platforme.

Rijndael je u tom poređenju imao veoma dobre softverske performanse na
tadašnjim 8-, 32- i 64-bitnim arhitekturama, nizak memorijski zahtev i
dobru pogodnost za implementacije na ograničenim uređajima
\cite{nechvatal2001aes,daemen2002rijndael}. NIST je istovremeno
Rijndaelovu bezbednosnu marginu smatrao adekvatnom, ali je dokumentovao i
kritike da je ona kod pojedinih varijanti manja nego kod nekih drugih
finalista. Važna karakteristika procesa je upravo to što takve primedbe
nisu skrivene: pretpostavljalo se da će kriptoanaliza napredovati i da će
se standard periodično preispitivati \cite{nechvatal2001aes}.

Iz današnje perspektive posebno je važno šta tadašnja dokumentacija
\emph{nije} predviđala. Efikasnost Rijndaela odnosila se na tada dostupne
softverske i hardverske platforme i na osobine dizajna koje su omogućavale
efikasnu realizaciju. Iz toga ne sledi da je AES proces predvideo buduće
namenske CPU instrukcije, poput AES-NI, niti GPU/SIMT izvršavanje,
savremene višestepene memorijske hijerarhije ili PCIe troškove.
Originalna procena može zato da podrži tvrdnju da je Rijndael izabran kao
algoritam sa dobrim implementacionim osobinama, ali ne i retroaktivnu
tvrdnju da su konkretni mehanizmi ubrzanja iz 2010-ih i 2020-ih bili deo
originalnog očekivanja.

Kontinuitet tri FIPS~197 varijante do ažuriranog izdanja potvrđuje
dugovečnost osnovne konstrukcije~\cite{nist_fips197_2001,nist_fips197_2023}.

\subsection{Migracija, zastupljenost i trajanje standarda}

Očekivana zamena DES-a ostvarena je njegovim povlačenjem, uz TDEA kao
prelaznu mogućnost~\cite{federalregister2005des}. TDEA je zatim opstao
znatno duže, uz postepena ograničenja i konačno povlačenje za novu
zaštitu~\cite{nist_sp80067r2,nist_sp800131ar2,nist_tdea_withdrawal_2023}.
Detaljan normativni sled dat je u odeljku~\ref{subsec:tdea_withdrawal}.
Retrospektivno, ta duga tranzicija pokazuje značaj instalirane baze,
interoperabilnosti i troška migracije: pojava naslednika i prestanak
upotrebe prethodnika nisu isti događaj.

Slična razlika je važna i za AES-128, AES-192 i AES-256. Sva tri su i
2026. deo FIPS~197 \cite{nist_fips197_2023}, ali to ne znači da imaju
jednaku zastupljenost u svim protokolima i sistemima. TLS~1.3, na primer,
zahteva implementaciju skupa \texttt{TLS\_AES\_128\_GCM\_SHA256}, dok
su \texttt{TLS\_AES\_256\_GCM\_SHA384} i ChaCha20-Poly1305 među
preporučenim dodatnim skupovima \cite{rfc8446}. Ovo je dokaz
standardizacione i protokolarne uloge AES-a, ali ne predstavlja samo po
sebi globalnu telemetriju o relativnom udelu AES-128, AES-192 i AES-256.

\subsection{Bezbednost, performanse i promena fokusa}

Više od dve decenije kriptoanalize značajno je produbilo razumevanje
AES-a, ali nije poništilo osnovnu bezbednosnu pretpostavku na kojoj je
standardizacija zasnovana. NIST je 2021. objavio formalni pregled
FIPS~197 i dosadašnje analize AES-a \cite{nist_ir8319}. Istovremeno,
akademska literatura sadrži važne rezultate nad punim brojem rundi.
Biryukov, Khovratovich i Nikolić pokazali su 2009. napad na puni
AES-256 u specijalizovanom related-key modelu
\cite{biryukov2009aes256}, dok su Bogdanov, Khovratovich i Rechberger
2011. predstavili biclique kriptoanalizu punog AES-a, sa složenošću
nešto manjom od iscrpne pretrage za sve tri dužine ključa
\cite{bogdanov2011biclique}. Ovi rezultati su kriptoanalitički značajni,
ali ne predstavljaju praktičan oporavak ključa u uobičajenom modelu
primene AES-a.

Retrospektivno je zato korisnije govoriti o promeni fokusa nego o tome da
je analiza primitive prestala. Dok je matematička analiza AES-a
nastavljena, praktična bezbednost sve više je zavisila od načina na koji
se algoritam implementira i uključuje u protokole. Implementacioni razvoj iz odeljka~\ref{subsec:aes_sidechannels}
pokazuje kako AES-NI povezuje performansni napredak sa uklanjanjem
važnog izvora cache/timing curenja~\cite{gueron2012aesni}, bez opšte
garancije protiv svih bočnih kanala.

Eksperimentalnu vezu sa ovom promenom pruža poređenje iz
odeljka~\ref{sec:aes-akceleracija}. Softverska i AES-NI putanja tamo
izvršavaju istu AES primitivu, dok razlika u vremenu pokazuje doprinos
kasnije uvedenih procesorskih instrukcija. Izmereni faktor ubrzanja zato
nije retroaktivna potvrda da je NIST 2001. predvideo AES-NI. On pokazuje
da se performanse istog standardizovanog algoritma mogu radikalno
promeniti promenom implementacione putanje.

Promena fokusa vidljiva je i na nivou režima rada. NIST je 2007.
standardizovao GCM kao režim autentifikovane enkripcije sa pridruženim
podacima, kombinujući poverljivost i autentifikaciju
\cite{nist_sp80038d}. TLS~1.3 zatim je uklonio ranije ne-AEAD skupove
šifara i definisao AEAD skupove, pri čemu je
\texttt{TLS\_AES\_128\_GCM\_SHA256} obavezan za implementaciju
\cite{rfc8446}. Savremeni IETF BCP dodatno ograničava upotrebu CBC
skupova u TLS~1.2 i preporučuje savremenije režime i verzije protokola
\cite{rfc9325}.

Ishod je širi predmet bezbednosne procene: uz AES konstrukciju,
protokoli moraju ispuniti zahteve načina rada, uključujući GCM pravila
o IV vrednostima~\cite{nist_sp80038d}.

\subsection{Kako se ostvaruju AES performanse: 2001--2003. i 2026.}

Tokom AES procesa performanse su razmatrane kao osobina algoritma na
različitim tadašnjim platformama. Rijndael je pokazivao dobre rezultate
u softveru, na procesorima različite širine reči i na ograničenim
uređajima \cite{nechvatal2001aes}. Efikasne softverske realizacije mogle
su kombinovati više operacija AES runde u unapred izračunate tabele,
čime se smanjivao broj aritmetičkih operacija po bloku
\cite{daemen2002rijndael}. Takav način izvršavanja predstavlja jedan
istorijski sloj performansi, ali ne i konačni oblik u kome se AES
izvršava na savremenim procesorima.

Arhitekturske optimizacije i namenske instrukcije čine sledeći sloj
~\cite{gueron2012aesni}. Savremene performanse tako spajaju originalni
dizajn sa kasnijim razvojem procesora i biblioteka.

Treći sloj je paralelno izvršavanje na GPU uređajima. CUDA/SIMT model
omogućava da veliki broj nezavisnih blokova bude obrađen istovremeno.
Literatura pokazuje da visoko optimizovani GPU AES kerneli mogu ostvariti
vrlo visoku propusnost, ali takvi rezultati zavise od implementacije,
memorijske hijerarhije i tačno definisanog obuhvata merenja
\cite{tezcan2021gpu}. Kernel-only rezultat nije isto što i propusnost
kompletne aplikacione putanje.

Ta razlika je direktno vidljiva u eksperimentu iz
odeljka~\ref{sec:cpu-gpu-poredjenje}. Odvojeni resident i end-to-end obuhvati daju različite ocene GPU prednosti. Rezultati pokazuju da CUDA kernel za veće
bafere može nadmašiti CPU AES-NI kada se posmatra izolovano računanje,
dok kompletna testirana pipeline putanja ne ostvaruje prednost nad
AES-NI. Time se dobija važna retrospektivna pouka: originalna osobina
algoritma da se može efikasno implementirati opstala je, ali konkretan
izvor performansi 2026. godine često leži izvan same primitive.

AES-GCM eksperiment dodatno pokazuje četvrti sloj. Savremena
kriptografska performansa nije samo brzina blokovske šifre, jer GCM uz
AES-CTR obradu uključuje i autentifikacionu komponentu. Zato CBC, čisti
AES kernel i AES-GCM ne predstavljaju iste operacije i njihove MB/s
vrednosti ne treba koristiti kao da mere isti posao. Ovo je u skladu sa
širim pomeranjem od pitanja ``koliko je brza blokovska šifra'' ka
pitanju ``koliko košta kompletna autentifikovana obrada podataka''.

\subsection{Dužina bloka i retrospektivna pitanja}

Jedna od trajnih razlika između DES/TDEA i AES familije je veličina
bloka. DES i TDEA rade sa blokovima od 64 bita, dok FIPS~197 definiše
128-bitne blokove za sve tri AES varijante
\cite{nist_fips197_2023}. Veličina bloka postaje važna kada se velika
količina podataka obrađuje istim ključem, jer se verovatnoća kolizije
povećava približavanjem birthday granici od reda $2^{n/2}$ blokova za
$n$-bitnu blokovsku šifru.

Praktičnu važnost tog problema pokazao je Sweet32. Za blok od 64 bita,
birthday granica je reda $2^{32}$ blokova, odnosno približno 32~GB
podataka, što je dovoljno malo da postane relevantno za dugotrajne
mrežne sesije. Bhargavan i Leurent su pokazali praktične kolizione
napade na 64-bitne blokovske šifre u TLS-u i OpenVPN-u
\cite{bhargavan2016sweet32}. AES sa blokom od 128 bita pomera
analogni generički birthday obim na približno $2^{64}$ blokova, pa je
specifični problem koji je pogodio TDEA pomeren daleko izvan uobičajenih
količina podataka.

To, međutim, ne znači da AES ima neograničen dozvoljeni obim pod jednim
ključem. Granice zavise od režima rada. SP~800-38D, na primer, ograničava
dužinu jednog GCM otvorenog teksta i postavlja posebne uslove za konstrukciju
i broj nonce/IV vrednosti; za slučajno generisane IV vrednosti broj
invokacija pod jednim ključem ograničen je na $2^{32}$
\cite{nist_sp80038d}. To je režimsko ograničenje, a ne ``lom'' AES
primitive.

Zanimljivo je da se ovo pitanje ponovo otvorilo upravo pred kraj
posmatranog perioda. NIST je 2024. pokrenuo proces razmatranja
Rijndael-256, varijante sa 256-bitnim blokom, navodeći da 128-bitni AES
blok ostaje dovoljan za mnoge primene, ali da veoma veliki obimi podataka
mogu opravdati širi blok \cite{nist_sp800197_predraft}. Tokom 2026.
NIST je u okviru revizije SP~800-38D zatražio komentare i o predloženom
wGCM režimu zasnovanom na blokovskoj šifri sa 256-bitnim blokom
\cite{nist_wgcm_2026}. Ovi predlozi nisu deo FIPS~197 AES standarda
2026. godine, ali pokazuju da pitanje dozvoljenog obima podataka ostaje
aktuelno čak i nakon četvrt veka primene AES-a.

Retrospektivni zaključak zato nije da je 128-bitni blok ``rešio'' problem
zauvek, već da je pružio znatno veći prostor od DES/TDEA blokova i
izbegao praktičnu Sweet32 klasu problema, dok savremeni režimi i
visokopropusni sistemi i dalje zahtevaju eksplicitno upravljanje
količinom podataka, nonce vrednostima i promenom ključa.

\subsection{Kvantna perspektiva bez senzacionalizma}

Kvantna analiza AES-a zahteva posebno pažljivo razdvajanje asimptotskog
rezultata od praktičnog napada. Groverov algoritam omogućava pretragu
neuređenog prostora veličine $N$ sa reda $\sqrt{N}$ kvantnih upita,
pa se za idealizovanu iscrpnu pretragu $k$-bitnog ključa često navodi
red veličine $2^{k/2}$ oracle poziva \cite{grover1996}. Iz toga potiče
popularna skraćenica da AES-128 ima ``64-bitnu kvantnu složenost'', ali
ona opisuje samo idealizovani broj upita i nije potpuna procena resursa
realnog kriptoanalitičkog sistema.

AES oracle mora biti realizovan kao reverzibilno kvantno kolo, uz
konkretan broj logičkih kubita, vrata i dubinu kola. Grassl i saradnici
su dali rane konkretne procene resursa za Groverovu pretragu AES ključa
\cite{grassl2016grover}, a kasniji rad Jaquesa i saradnika dodatno je
analizirao konkretan trošak AES oracle-a pod ograničenjima dubine
\cite{jaques2020grover}. Takve procene pokazuju zašto samo zamena
$2^k$ sa $2^{k/2}$ nije dovoljna za procenu praktičnog napada: relevantni
su serijska dubina, broj kvantnih vrata, broj logičkih i fizičkih kubita,
korekcija grešaka i mogućnost paralelizacije.

Aktuelna NIST interpretacija dodatno ublažava senzacionalističko
tumačenje Grovera. NIST navodi da puna kvadratna prednost zahteva veliki
broj serijskih koraka i da masovna paralelizacija, neophodna za praktične
napade, smanjuje relativnu kvantnu prednost. NIST zato ne preporučuje
automatsko udvostručavanje simetričkih ključeva samo zbog Groverovog
algoritma i eksplicitno navodi da trenutne aplikacije mogu nastaviti da
koriste AES sa ključevima od 128, 192 ili 256 bita
\cite{nist_pqc_faq}.

AES-256 ipak pruža veći prostor ključeva i time veću rezervu i u
klasičnom i u idealizovanom kvantnom modelu. To nije isto što i tvrdnja
da je AES-128 danas nebezbedan, niti da AES-256 ``rešava'' sve
postkvantne probleme. Najveći uticaj poznatih kvantnih algoritama na
savremenu kriptografsku tranziciju odnosi se na asimetrične primitive
pogođene Shorovim algoritmom, dok za AES nema odgovarajućeg eksponencijalnog
kvantnog ubrzanja. Kvantna perspektiva zato predstavlja dodatni faktor u
izboru bezbednosne margine, a ne dokaz zastarelosti AES-a.

Tabela~\ref{tab:retrospektiva} objedinjuje očekivanja i ishode, uz
razdvajanje potvrđenih kriterijuma od kasnijih tehnoloških promena.

\begin{table}[htbp]
\centering
\small
\setlength{\tabcolsep}{4pt}
\caption{Dokumentovana očekivanja iz perioda standardizacije AES-a i
proverljivi ishodi do 2026. godine. Ocena je opisna i ne predstavlja
numeričko rangiranje.}
\label{tab:retrospektiva}
\begin{tabular}{@{}>{\raggedright\arraybackslash}p{0.15\textwidth} >{\raggedright\arraybackslash}p{0.25\textwidth} >{\raggedright\arraybackslash}p{0.35\textwidth} >{\raggedright\arraybackslash}p{0.15\textwidth}@{}}
\toprule
Tema &
Dokumentovano očekivanje / kriterijum &
Ishod do 2026. &
Retrospektivna ocena \\
\midrule
Bezbednost primitive &
AES proces daje najveću težinu bezbednosti i procenjuje da Rijndael ima
adekvatnu marginu, uz očekivanje daljeg napretka kriptoanalize
\cite{nechvatal2001aes}. &
Related-key i biclique radovi proširili su teorijsku analizu punog AES-a,
ali nisu doveli do praktičnog oporavka ključa u standardnom modelu primene
\cite{biryukov2009aes256,bogdanov2011biclique,nist_ir8319}. &
U skladu sa očekivanjem, uz ogradu. \\

\addlinespace
Softverska efikasnost &
Rijndael je ocenjen kao veoma efikasan na različitim tadašnjim
softverskim platformama i ograničenim uređajima
\cite{nechvatal2001aes}. &
Savremeni softverski AES koristi arhitekturske optimizacije i često
namenske instrukcije; brzina 2026. nije posledica samo generičke
softverske implementacije. &
Delimično u skladu; kontekst promenjen. \\

\addlinespace
Hardverska efikasnost &
Hardverska realizacija i mogućnost paralelnog/pipeline izvršavanja bili
su deo kriterijuma AES procesa \cite{nechvatal2001aes}. &
AES je dobio namensku CPU podršku poput AES-NI i efikasne GPU
implementacije \cite{gueron2012aesni,tezcan2021gpu}. &
U skladu na opštem nivou; mehanizmi su kasniji. \\

\addlinespace
Namenske AES instrukcije &
AES-NI i slične buduće CPU instrukcije nisu bile dokumentovano
predviđanje AES procesa. &
Intel je od Westmere generacije uveo šest AES instrukcija sa
performansnim i implementaciono-bezbednosnim prednostima
\cite{gueron2012aesni}. &
Izvan originalnog obuhvata očekivanja. \\

\addlinespace
Dužine ključa &
Standard mora podržati 128-, 192- i 256-bitne ključeve
\cite{nechvatal2001aes,nist_fips197_2001}. &
Sve tri varijante ostaju deo FIPS~197 i posle ažuriranja iz 2023.
\cite{nist_fips197_2023}. &
U skladu sa kasnijim ishodom. \\

\addlinespace
DES migracija &
AES je razvijen kao naslednik zastarelog DES-a. &
FIPS~46-3 povučen je 19.\ maja 2005, uz preporuku prelaska na AES ili,
za prelazne primene, TDEA \cite{federalregister2005des}. &
U skladu sa kasnijim ishodom. \\

\addlinespace
TDEA kao prelazno rešenje &
TDEA je zadržan kao odobrena prelazna alternativa nakon pojave AES-a. &
Nova primena TDEA zaštite nedozvoljena je posle 31.\ decembra 2023, a
SP~800-67 Rev.~2 povučen je 1.\ januara 2024.
\cite{nist_sp800131ar2,nist_tdea_withdrawal_2023}. &
Duga tranzicija; konačno povučen. \\
\bottomrule
\end{tabular}
\end{table}

\begin{table}[htbp]
\centering
\small
\setlength{\tabcolsep}{4pt}
\ContinuedFloat
\caption[]{Dokumentovana očekivanja i ishodi do 2026. godine (nastavak).}
\begin{tabular}{@{}>{\raggedright\arraybackslash}p{0.15\textwidth} >{\raggedright\arraybackslash}p{0.25\textwidth} >{\raggedright\arraybackslash}p{0.35\textwidth} >{\raggedright\arraybackslash}p{0.15\textwidth}@{}}
\toprule
Tema &
Dokumentovano očekivanje / kriterijum &
Ishod do 2026. &
Retrospektivna ocena \\
\midrule
Veličina bloka &
AES standardizuje blok od 128 bita, nasuprot 64-bitnom DES/TDEA bloku
\cite{nist_fips197_2001}. &
Sweet32 je praktično pokazao ograničenja 64-bitnog bloka pri velikom
obimu podataka; NIST 2024--2026 istražuje i šire 256-bitne blokove za
buduće visokopropusne primene
\cite{bhargavan2016sweet32,nist_sp800197_predraft,nist_wgcm_2026}. &
128 bita je dobro ostarilo, ali nije bez granica. \\

\addlinespace
Implementa\-ciona bezbednost &
AES proces je razmatrao implementacione karakteristike i bočne kanale,
ali u kontekstu tadašnjeg hardvera \cite{nechvatal2001aes}. &
Cache/timing rizici tabelarnih implementacija postali su važan praktičan
problem; AES-NI uklanja lookup tabele i podatkovno zavisno vreme za tu
klasu implementacije \cite{gueron2012aesni}. &
Promenjen implementacioni kontekst. \\

\addlinespace
AEAD i GCM &
GCM nije bio deo originalnog izbora AES primitive; režimi rada
standardizovani su odvojeno. &
SP~800-38D definiše GCM kao AEAD, a TLS~1.3 zahteva implementaciju
AES-128-GCM skupa \cite{nist_sp80038d,rfc8446}. &
Izvan originalnog obuhvata; protokolarni razvoj. \\

\addlinespace
GPU/SIMT izvršavanje &
Savremeni GPU compute model nije bio dokumentovano očekivanje AES procesa. &
AES ima efikasne GPU realizacije, ali kernel propusnost i kompletna
host--device putanja predstavljaju različite metrike
\cite{tezcan2021gpu}. &
Izvan originalnog obuhvata. \\

\addlinespace
Kvantna pretraga &
AES proces je primarno procenjivao klasičnu iscrpnu pretragu i poznatu
klasičnu kriptoanalizu. &
Grover daje kvadratno ubrzanje u idealizovanom modelu, ali konkretni
resursi, serijska dubina i paralelizacija znatno komplikuju praktičnu
procenu \cite{grassl2016grover,jaques2020grover,nist_pqc_faq}. &
Izvan originalnog obuhvata; zahteva ogradu. \\
\bottomrule
\end{tabular}
\end{table}

\FloatBarrier

